\section{Introduction}

Teleoperation brings human judgment to remote manipulation tasks that remain difficult for autonomous robots, including hazardous operation, flexible manufacturing, and unstructured manipulation. Haptic feedback can convey remote contact information to the operator and support assessment of contact state, grasp stability, and slip. Consequently, a practical teleoperation system must integrate bilateral control with perception, mechanical execution, and the human--machine interface \cite{ref01,ref02}.

For heterogeneous objects, one fixed setting cannot simultaneously provide gentle contact and robust positioning, grasping, and transport. Lower stiffness and grasp force can reduce contact severity but can impair positioning and holding stability; higher settings can improve positional robustness while increasing the risk of impact or deformation. This trade-off motivates object-conditioned configuration across the subsystems involved in telemanipulation, rather than a single fixed setting for all objects.

Impedance control regulates the relationship among displacement, velocity, and interaction force and is a standard means of achieving compliant robot behavior \cite{ref03}. Variable-impedance methods adapt stiffness and damping to task state, contact information, learned policies, or human state, while addressing the stability implications of time-varying parameters \cite{ref04,ref05}. In teleoperation, tele-impedance and bilateral-control studies have regulated remote-robot stiffness according to operator motion, estimated human impedance, contact state, or task phase \cite{ref06,ref07,ref08,ref09,ref18,ref19,ref20}. Related work has also examined operator-facing haptic feedback and tactile feedback for telemanipulation \cite{ref17,ref22,ref25}. These studies establish important approaches to remote mechanical compliance and haptic interaction, but their principal focus is commonly on impedance transmission, stiffness regulation, transparency, or feedback presentation.

Visual and multimodal information has extended semi-autonomous tele-impedance beyond contact-driven adaptation. Vision and voice have been used to select slave-side impedance from object properties with operator confirmation or correction \cite{ref10}; visual-force control has combined visual and force information in a visual-feature space \cite{ref11}; and task-oriented shared control has reduced manual decision-making \cite{ref12}. More recent work has derived stiffness from object geometry, material, and environmental relationships \cite{ref13}, or generated task-relevant stiffness matrices from visual, gaze, and language inputs \cite{ref14}. Automatic switching and passive arbitration address when control authority or force/stiffness scaling should change \cite{ref15,ref16}, while teleoperated grasping has used tactile deformation and grip-safety information to regulate gripper impedance \cite{ref30}. Thus, object- and task-dependent parameter adaptation is established; in particular, visual information can be used to select or generate remote impedance settings.

The remaining system-level question is not whether vision can adjust impedance, but whether one interpretable object-conditioned strategy can jointly configure the remote arm, the operator-facing haptic interface, and gripper execution. In the closely related literature considered here, these elements are usually adjusted as separate design components or the object information is translated primarily into remote stiffness, impedance matrices, or control-authority decisions. We therefore investigate whether a common parameter bundle spanning these three subsystems provides a task-level benefit beyond visual-semantic display, manual strategy selection, or impedance-only reconfiguration.

Joint configuration also creates an integration problem: visual inference is intermittent and can be uncertain, whereas the teleoperation update must remain responsive to the operator. Passing every perception result directly into the control path can either delay the supervisory update or induce repeated switching when detections fluctuate. The required architecture must therefore transfer semantic information into a limited and consistent configuration action without turning the system into continuous visual servoing.

To address this question, we propose an object-conditioned joint-configuration framework for haptic teleoperation of heterogeneous objects. Each detected object class is mapped to one of three operational strategies and then to a parameter bundle specifying slave-side impedance, master-side haptic rendering, and gripper execution. After task initiation, the first threshold-qualified mapped detection can initiate a one-shot configuration event, after which the selected strategy is locked for the remainder of the trial. A multi-rate implementation keeps visual inference outside the supervisory teleoperation call path and prevents repeated reconfiguration caused by fluctuating detections. The framework is evaluated on a closed set of six known object instances using five modes: a fixed-parameter baseline (A), a manual-selection workflow baseline (B), joint parameter-bundle configuration (C), visual-semantic display with fixed parameters (D), and impedance-only reconfiguration (E). The primary C--E comparison tests the system-level contribution of the complete bundle under matched visual-semantic and impedance conditions, whereas C--D tests the contribution beyond the displayed visual-semantic information.

The main contributions are as follows:

\begin{enumerate}
\item \textbf{Object-conditioned cross-subsystem configuration.} An interpretable strategy layer maps each detected object class to an operational strategy and jointly assigns a parameter bundle across slave-side impedance, master-side haptic rendering, and gripper execution.
\item \textbf{Multi-rate one-shot integration.} A non-blocking architecture decouples low-rate visual inference from the teleoperation update and converts a qualified semantic result into a bounded, one-shot configuration event with intra-trial strategy locking.
\item \textbf{Matched human-in-the-loop evaluation.} A five-mode experiment on an Omega.7--Panda platform, involving three operators, six heterogeneous objects, and 135 trials, evaluates whether joint configuration provides an additional system-level benefit over visual-semantic display and impedance-only reconfiguration.
\end{enumerate}

The evaluation is designed to assess the bundled configuration at the system level. It does not identify the independent causal contributions of the haptic-interface and gripper parameter groups, and its empirical scope is limited to the six known objects and repeated measurements in the present sample.
